% !TEX root = ../paper_zh.tex

\section{结论}

本文提出 UAV-TAlign-12K，一个基准规模的无人机红外-可见光对齐资源。它包含 6,039 对同步
图像、经过完整性检查的 6,037 对正式评测划分，以及覆盖成对证据、场景共识和可靠性--覆盖率
分析的多层协议。七种代表性基线表明，红外专用方法和现代匹配器已经能够在真实采集跨视场条件下
提供高可用性的成对单应性。

UAV-TAlign 将这些证据转化为 15 个场景变换。框架结合渐进式调度、几何校准帧加权、鲁棒共识
和跨模态验证，并在规范运行点识别出 9 个高置信产品，覆盖 74.4\% 的评测图像对。留帧重采样和
受控扰动进一步支撑了场景诊断量的稳定性及其对几何变化的敏感性。这些结果表明，场景级证据整合
与置信度感知运行分析是跨视场无人机红外-可见光对齐中的核心评测维度。

\section*{数据可用性}

UAV-TAlign-12K、正式评测 manifest、元数据、完整性报告和评测脚本均已为公开发布做好准备。
正式评测划分包含 6,039 对无人机红外-可见光对齐数据中的 6,037 对完整性检查 RGB-红外图像。
数据集归广西大学所有，将采用知识共享署名 4.0 国际许可协议发布；代码与评测材料将通过公共
仓库提供。同行评审期间可提供审稿访问权限，正式 manifest 哈希将用于保证评测可复现性。
